% GENERATED FILE. Edit canonical lessons, then run npm run generate:books.
% canonical-source-hash: fnv1a64:0a171e62bff1a101
% canonical-lessons: IT-C422-battello, IT-C422-gommone, IT-C422-remo, IT-C422-remare, IT-C422-canoa

\chapter{In barca}
\label{ch:it-in-barca}
\hlchaptermodality{422}

\begin{chapteropening}
\textbf{By the end of this chapter:} \emph{I can talk about small boats: the lake boat, the dinghy, the oar, rowing and the canoe.}

Complete the last lesson of chapter 422: I can talk about small boats: the lake boat, the dinghy, the oar, rowing and the canoe.
\end{chapteropening}

\section[il battello]{il battello — the boat (on a lake or river), the small steamer}

\label{lesson:IT-C422-battello}



\begin{quote}
Before the new one: say the Italian for to disembark, then the Italian for to set sail.
\end{quote}

\subsection*{You'll want to know: il battello}
\textbf{il battello} — \textquotedblleft{}the boat (on a lake or river), the small steamer\textquotedblright{}.

\begin{quote}
Sul lago prendiamo il battello.
\end{quote}

\begin{quote}
\emph{On the lake we take the boat.}
\end{quote}

For lakes and rivers; on the sea you take \emph{il traghetto}.

\begin{cousinweb}
\emph{battello} was borrowed from Old French \emph{batel}, a little form of a word taken from Old English \emph{bāt} 'boat' (EtymDB).
\end{cousinweb}

\begin{grammarlens}[title={What you've built}]
One more word for telling, explaining and describing at B1.
\end{grammarlens}

\subsection*{Your turn}
\begin{itemize}
\raggedright
  \item \emph{Say it:} \emph{il battello}
  \item \emph{Say it:} \emph{il battello}, once more
  \item \emph{Say it:} a sentence of your own with \emph{il battello}
  \item \emph{Recall:} say the Italian for to disembark, then the Italian for to set sail, then say \emph{il battello} again
\end{itemize}

\begin{tcolorbox}[breakable,colback=teal!4,colframe=teal!35!black,title={Before you move on}]
What does il battello mean? (\textquotedblleft{}The boat (on a lake or river), the small steamer\textquotedblright{}.) Say it once more.
\end{tcolorbox}
\section[il gommone]{il gommone — the rubber dinghy, the inflatable boat}

\label{lesson:IT-C422-gommone}



\begin{quote}
Before the new one: say the Italian for to set sail, then the Italian for the boat.
\end{quote}

\subsection*{You'll want to know: il gommone}
\textbf{il gommone} — \textquotedblleft{}the rubber dinghy, the inflatable boat\textquotedblright{}.

\begin{quote}
Con il gommone andiamo alle grotte.
\end{quote}

\begin{quote}
\emph{With the dinghy we go to the caves.}
\end{quote}

A big thing made of \emph{gomma}, rubber.

\begin{cousinweb}
\emph{gommone} is built on \emph{gomma} 'rubber', which is inherited from Late Latin \emph{gumma}, from Latin \emph{gummi} 'gum' (EtymDB).
\end{cousinweb}

\begin{grammarlens}[title={What you've built}]
One more word for telling, explaining and describing at B1.
\end{grammarlens}

\subsection*{Your turn}
\begin{itemize}
\raggedright
  \item \emph{Say it:} \emph{il gommone}
  \item \emph{Say it:} \emph{il gommone}, once more
  \item \emph{Say it:} a sentence of your own with \emph{il gommone}
  \item \emph{Recall:} say the Italian for to set sail, then the Italian for the boat, then say \emph{il gommone} again
\end{itemize}

\begin{tcolorbox}[breakable,colback=teal!4,colframe=teal!35!black,title={Before you move on}]
What does il gommone mean? (\textquotedblleft{}The rubber dinghy, the inflatable boat\textquotedblright{}.) Say it once more.
\end{tcolorbox}
\section[il remo]{il remo — the oar}

\label{lesson:IT-C422-remo}



\begin{quote}
Before the new one: say the Italian for the boat, then the Italian for the rubber dinghy.
\end{quote}

\subsection*{You'll want to know: il remo}
\textbf{il remo} — \textquotedblleft{}the oar\textquotedblright{}.

\begin{quote}
Ho perso un remo nell'acqua.
\end{quote}

\begin{quote}
\emph{I lost an oar in the water.}
\end{quote}

\emph{Una barca a remi} — a rowing boat.

\begin{cousinweb}
\emph{remo} is inherited from Latin \emph{rēmus} 'oar' (EtymDB).
\end{cousinweb}

\begin{grammarlens}[title={What you've built}]
One more word for telling, explaining and describing at B1.
\end{grammarlens}

\subsection*{Your turn}
\begin{itemize}
\raggedright
  \item \emph{Say it:} \emph{il remo}
  \item \emph{Say it:} \emph{il remo}, once more
  \item \emph{Say it:} a sentence of your own with \emph{il remo}
  \item \emph{Recall:} say the Italian for the boat, then the Italian for the rubber dinghy, then say \emph{il remo} again
\end{itemize}

\begin{tcolorbox}[breakable,colback=teal!4,colframe=teal!35!black,title={Before you move on}]
What does il remo mean? (\textquotedblleft{}The oar\textquotedblright{}.) Say it once more.
\end{tcolorbox}
\section[remare]{remare — to row}

\label{lesson:IT-C422-remare}



\begin{quote}
Before the new one: say the Italian for the rubber dinghy, then the Italian for the oar.
\end{quote}

\subsection*{You'll want to know: remare}
\textbf{remare} — \textquotedblleft{}to row\textquotedblright{}.

\begin{quote}
Abbiamo remato fino all'altra riva.
\end{quote}

\begin{quote}
\emph{We rowed to the other bank.}
\end{quote}

Figuratively, \emph{remare contro} — to work against someone.

\begin{cousinweb}
\emph{remare} is built on \emph{remo} 'oar' (EtymDB), which is inherited from Latin \emph{rēmus} 'oar'.
\end{cousinweb}

\begin{grammarlens}[title={What you've built}]
One more word for telling, explaining and describing at B1.
\end{grammarlens}

\subsection*{Your turn}
\begin{itemize}
\raggedright
  \item \emph{Say it:} \emph{remare}
  \item \emph{Say it:} \emph{remare}, once more
  \item \emph{Say it:} a sentence of your own with \emph{remare}
  \item \emph{Recall:} say the Italian for the rubber dinghy, then the Italian for the oar, then say \emph{remare} again
\end{itemize}

\begin{tcolorbox}[breakable,colback=teal!4,colframe=teal!35!black,title={Before you move on}]
What does remare mean? (\textquotedblleft{}To row\textquotedblright{}.) Say it once more.
\end{tcolorbox}
\section[la canoa]{la canoa — the canoe}

\label{lesson:IT-C422-canoa}



\begin{quote}
Before the new one: say the Italian for the oar, then the Italian for to row.
\end{quote}

\subsection*{You'll want to know: la canoa}
\textbf{la canoa} — \textquotedblleft{}the canoe\textquotedblright{}.

\begin{quote}
D'estate facciamo canoa sul fiume.
\end{quote}

\begin{quote}
\emph{In summer we go canoeing on the river.}
\end{quote}

\emph{Andare in canoa} — to go canoeing.

\begin{cousinweb}
\emph{canoa} was borrowed from Spanish \emph{canoa}, which took it from Taíno \emph{kanoa} 'dugout canoe' (EtymDB), a language of the Caribbean.
\end{cousinweb}

\begin{grammarlens}[title={What you've built}]
One more word for telling, explaining and describing at B1.
\end{grammarlens}

\subsection*{Your turn}
\begin{itemize}
\raggedright
  \item \emph{Say it:} \emph{la canoa}
  \item \emph{Say it:} \emph{la canoa}, once more
  \item \emph{Say it:} a sentence of your own with \emph{la canoa}
  \item \emph{Recall:} say the Italian for the oar, then the Italian for to row, then say \emph{la canoa} again
\end{itemize}

\begin{tcolorbox}[breakable,colback=teal!4,colframe=teal!35!black,title={Before you move on}]
What does la canoa mean? (\textquotedblleft{}The canoe\textquotedblright{}.) Say it once more.
\end{tcolorbox}
